% problems5.tex — Problem Set and Answers for Chapter 5

\begin{problemset}

\subsection*{Analytical Exercises}

\exercise[1]
(Fixed-effects estimator as the LSDV estimator)\quad
In \eqref{eq:5.1.1pp} on page~329, redefine $\alpha_i$ to be the sum of old $\alpha_i$ and
$\mathbf{b}_i'\boldsymbol{\gamma}$, so that \eqref{eq:5.1.1pp} becomes
\begin{equation}
  \mathbf{y}_i = \mathbf{F}_i\boldsymbol{\beta} + \mathbf{1}_M \cdot \alpha_i + \boldsymbol{\eta}_i
  \quad (i = 1, 2, \ldots, n). \tag{1}
\end{equation}

Define the following stacked vectors and matrices:
\[
  \underset{(Mn \times 1)}{\mathbf{y}} = \begin{bmatrix} \mathbf{y}_1 \\ \vdots \\ \mathbf{y}_n \end{bmatrix}, \quad
  \underset{(Mn \times \#\boldsymbol{\beta})}{\mathbf{F}} = \begin{bmatrix} \mathbf{F}_1 \\ \vdots \\ \mathbf{F}_n \end{bmatrix}, \quad
  \underset{(Mn \times 1)}{\boldsymbol{\eta}} = \begin{bmatrix} \boldsymbol{\eta}_1 \\ \vdots \\ \boldsymbol{\eta}_n \end{bmatrix}, \quad
  \underset{(n \times 1)}{\boldsymbol{\alpha}} = \begin{bmatrix} \alpha_1 \\ \vdots \\ \alpha_n \end{bmatrix},
\]
\[
  \underset{(Mn \times (n + \#\boldsymbol{\beta}))}{\mathbf{W}} = (\mathbf{D} \vdots \mathbf{F}), \quad
  \underset{((n + \#\boldsymbol{\beta}) \times 1)}{\boldsymbol{\theta}}
  = \begin{bmatrix} \boldsymbol{\alpha} \\ \boldsymbol{\beta} \end{bmatrix},
\]
\begin{equation}
  \underset{(Mn \times n)}{\mathbf{D}} = \mathbf{I}_n \otimes \mathbf{1}_M
  = \begin{bmatrix}
    \mathbf{1}_M & & \\
    & \mathbf{1}_M & \\
    & & \ddots \\
    & & & \mathbf{1}_M
  \end{bmatrix}, \tag{2}
\end{equation}
so that the $M$-equation system (1) for $i = 1, 2, \ldots, n$ can be written as
\[
  \begin{bmatrix} \mathbf{y}_1 \\ \mathbf{y}_2 \\ \vdots \\ \mathbf{y}_n \end{bmatrix}
  = \begin{bmatrix} \mathbf{1}_M \cdot \alpha_1 \\ \mathbf{1}_M \cdot \alpha_2 \\ \vdots \\ \mathbf{1}_M \cdot \alpha_n \end{bmatrix}
  + \begin{bmatrix} \mathbf{F}_1 \\ \mathbf{F}_2 \\ \vdots \\ \mathbf{F}_n \end{bmatrix}\boldsymbol{\beta}
  + \begin{bmatrix} \boldsymbol{\eta}_1 \\ \boldsymbol{\eta}_2 \\ \vdots \\ \boldsymbol{\eta}_n \end{bmatrix}
\]
or
\begin{equation}
  \mathbf{y} = \mathbf{D}\boldsymbol{\alpha} + \mathbf{F}\boldsymbol{\beta} + \boldsymbol{\eta}
  = \mathbf{W}\boldsymbol{\theta} + \boldsymbol{\eta}. \tag{3}
\end{equation}

The OLS estimator of $\boldsymbol{\theta}$ on a pooled sample of size $Mn$, $(\mathbf{y}, \mathbf{W})$, is
$(\mathbf{W}'\mathbf{W})^{-1}\cdot\mathbf{W}'\mathbf{y}$.

\begin{enumerate}[label=(\alph*)]
\item Show that the fixed-effects estimator given in \eqref{eq:5.2.4} is the last $\#\boldsymbol{\beta}$ elements
  of this vector. \textbf{Hint:} Use the Frisch-Waugh Theorem about partitioned
  regressions (see Analytical Exercise~4(c) of Chapter~1). The first
  step is to regress $\mathbf{y}$ on $\mathbf{D}$ and regress each column of $\mathbf{F}$ on $\mathbf{D}$. If we define
  the annihilator associated with $\mathbf{D}$ as
  \begin{equation}
    \underset{(Mn \times Mn)}{\mathbf{M}_\mathbf{D}}
    \equiv \mathbf{I}_{Mn} - \mathbf{D}(\mathbf{D}'\mathbf{D})^{-1}\mathbf{D}', \tag{4}
  \end{equation}
  then the residual vectors can be written as $\mathbf{M}_\mathbf{D}\mathbf{y}$ and
  $\mathbf{M}_\mathbf{D}\mathbf{F}$. The second step
  is to regress $\mathbf{M}_\mathbf{D}\mathbf{y}$ on $\mathbf{M}_\mathbf{D}\mathbf{F}$.
  Show that the estimated regression coefficient
  can be written as
  \begin{equation}
    (\mathbf{F}'\mathbf{M}_\mathbf{D}\mathbf{F})^{-1}\mathbf{F}'\mathbf{M}_\mathbf{D}\mathbf{y}. \tag{5}
  \end{equation}
  Show that $\mathbf{M}_\mathbf{D} = \mathbf{I}_n \otimes \mathbf{Q}$ (where $\mathbf{Q}$ is the annihilator associated with $\mathbf{1}_M$, as
  in the text) and use this to show that (1) is indeed the fixed-effects estimator.

  Alternatively, you can go all the way back to the normal equations. If $\mathbf{a}$
  and $\mathbf{b}$ are the OLS estimators of $\boldsymbol{\alpha}$ and $\boldsymbol{\beta}$,
  then $(\mathbf{a}, \mathbf{b})$ satisfies the system
  of $n + \#\boldsymbol{\beta}$ linear simultaneous equations:
  \begin{equation}
    \begin{bmatrix} \mathbf{D}'\mathbf{D} & \mathbf{D}'\mathbf{F} \\
      \mathbf{F}'\mathbf{D} & \mathbf{F}'\mathbf{F} \end{bmatrix}
    \begin{bmatrix} \mathbf{a} \\ \mathbf{b} \end{bmatrix}
    = \begin{bmatrix} \mathbf{D}'\mathbf{y} \\ \mathbf{F}'\mathbf{y} \end{bmatrix}. \tag{6}
  \end{equation}
  Take the first $n$ equations and solve for $\mathbf{a}$ taking $\mathbf{b}$ as given. This should
  yield: $\mathbf{a} = (\mathbf{D}'\mathbf{D})^{-1}(\mathbf{D}'\mathbf{y} - \mathbf{D}'\mathbf{F}\mathbf{b})$.
  Then substitute this into the rest of the
  equations to solve for $\mathbf{b}$. (You can accomplish the same thing by using the
  formula for inverting partitioned matrices.)

\item Show:
  \[
    a_i = \bar{y}_i - \biggl(\frac{1}{M}\sum_{m=1}^{M}\mathbf{f}_{im}'\biggr)
    \hat{\boldsymbol{\beta}}_{\text{FE}},
  \]
  where $a_i$ is the OLS estimator of $\alpha_i$, $\mathbf{f}_{im}'$ is the $m$-th row of $\mathbf{F}_i$, and
  $\bar{y}_i = \frac{1}{M}(y_{i1} + \cdots + y_{iM})$.
  \textbf{Hint:} Use the normal equations for $\boldsymbol{\alpha}$.

\item Assume:
  \begin{enumerate}[label=(\roman*)]
  \item $\{\mathbf{y}_i, \mathbf{F}_i\}$ is i.i.d.,
  \item $\E(\boldsymbol{\eta}_i \mid \mathbf{F}_i) = \mathbf{0}$
    (which is stronger than the orthogonality conditions \eqref{eq:5.1.8b}),
  \item $\E(\boldsymbol{\eta}_i\boldsymbol{\eta}_i' \mid \mathbf{F}_i)
    = \sigma_\eta^2\,\mathbf{I}_M$ (the spherical error assumption), and
  \item $\mathbf{W}$ is of full column rank.
  \end{enumerate}
  Show that the assumptions of the Classical Regression Model are satisfied
  for the pooled regression (3), that is,
  \[
    \E(\boldsymbol{\eta} \mid \mathbf{W}) = \underset{(Mn \times 1)}{\mathbf{0}}.
  \]
  Also show that $\E(\boldsymbol{\eta}\boldsymbol{\eta}' \mid \mathbf{W}) = \sigma_\eta^2\,\mathbf{I}_{Mn}$.
  Develop the \emph{small sample} theory for
  $\mathbf{a}$ and $\mathbf{b}$. For example, are they unbiased? Let $SSR$ be the sum of squared
  residuals from the pooled regression (3). Is this $SSR$ the same as the $SSR$
  from the ``within'' regression \eqref{eq:5.2.15}?
\end{enumerate}

\exercise[2]
(GMM interpretation of FE)\quad
It is fairly well known for the case of two equations
($M = 2$) that the fixed-effects estimator is the OLS estimator in the
regression of first differences: $y_{i2} - y_{i1}$ on $\mathbf{z}_{i2} - \mathbf{z}_{i1}$.
This question generalizes
this argument to the case of arbitrary $M$.

We wish to write the fixed-effects estimator \eqref{eq:5.2.4} in the GMM format
\begin{equation}
  (\mathbf{S}_{\mathbf{xz}}'\hat{\mathbf{W}}\mathbf{S}_{\mathbf{xz}})^{-1}
  \mathbf{S}_{\mathbf{xz}}'\hat{\mathbf{W}}\mathbf{s}_{\mathbf{xy}} \tag{7}
\end{equation}
for properly defined $\mathbf{S}_{\mathbf{xz}}$, $\hat{\mathbf{W}}$, and $\mathbf{s}_{\mathbf{xy}}$.
Let $\mathbf{C}$ be a matrix such that
\begin{equation}
  \text{$\mathbf{C}$ is an $M \times (M-1)$ matrix of full column rank satisfying }
  \mathbf{C}'\mathbf{1}_M = \mathbf{0}. \tag{8}
\end{equation}

A prominent example of such a matrix $\mathbf{C}$ is the matrix of first differences:
\[
  \underset{(M \times (M-1))}{\mathbf{C}} = \begin{bmatrix}
    -1 & 0 & 0 & 0 & \cdots & 0 \\
    1 & -1 & 0 & 0 & \cdots & 0 \\
    0 & 1 & -1 & 0 & \cdots & 0 \\
    & & \ddots & \ddots & & \\
    0 & \cdots & 0 & 1 & -1 & 0 \\
    0 & \cdots & 0 & 0 & 1 & -1 \\
    0 & \cdots & 0 & 0 & 0 & 1
  \end{bmatrix}.
\]

Another example is an $M \times (M-1)$ matrix created by dropping one column,
say the last, from $\mathbf{Q}$, the annihilator associated with $\mathbf{1}_M$.

\begin{enumerate}[label=(\alph*)]
\item Verify that these two matrices satisfy (8).

\item For any given choice of $\mathbf{C}$, verify that you can derive the following system
  of $M - 1$ transformed equations by multiplying both sides of the system of
  $M$ equations \eqref{eq:5.1.1pp} on page~329 from left by $\mathbf{C}'$:
  \begin{equation}
    \hat{\mathbf{y}}_i = \hat{\mathbf{F}}_i\boldsymbol{\beta} + \hat{\boldsymbol{\eta}}_i \tag{9}
  \end{equation}
  where
  \[
    \underset{((M-1) \times 1)}{\hat{\mathbf{y}}_i} \equiv \mathbf{C}'\mathbf{y}_i, \quad
    \underset{((M-1) \times \#\boldsymbol{\beta})}{\hat{\mathbf{F}}_i} \equiv \mathbf{C}'\mathbf{F}_i, \quad
    \underset{((M-1) \times 1)}{\hat{\boldsymbol{\eta}}_i} \equiv \mathbf{C}'\boldsymbol{\eta}_i.
  \]

\item (optional)\quad If $\{\mathbf{y}_i, \mathbf{Z}_i\}$ satisfies the assumptions of the error-components model
  (which are \eqref{eq:5.1.1pp}, \eqref{eq:5.1.2}, \eqref{eq:5.1.8},
  \eqref{eq:5.1.4}, \eqref{eq:5.1.5}, and \eqref{eq:5.1.6} or
  \eqref{eq:5.1.6p} on page~324), then
  $\{\hat{\mathbf{y}}_i, \hat{\mathbf{F}}_i\}$ satisfies the assumptions of the
  random-effects model, i.e., for the $M-1$ equation system (9),

  (random sample) $\{\hat{\mathbf{y}}_i, \hat{\mathbf{F}}_i\}$ is i.i.d.,

  (orthogonality conditions) $\E(\hat{\boldsymbol{\eta}}_i \otimes \mathbf{x}_i) = \mathbf{0}$ with
  $\mathbf{x}_i = \text{union of}$
  $(\mathbf{z}_{i1}, \ldots, \mathbf{z}_{iM}) = \text{unique elements of } (\mathbf{F}_i, \mathbf{b}_i)$,

  (identification) $\E(\hat{\mathbf{F}}_i \otimes \mathbf{x}_i)$ is of full column rank,

  (conditional homoskedasticity) $\E(\hat{\boldsymbol{\eta}}_i\hat{\boldsymbol{\eta}}_i' \mid \mathbf{x}_i)
  = \E(\hat{\boldsymbol{\eta}}_i\hat{\boldsymbol{\eta}}_i')$ which is nonsingular,

  (nonsingularity of $\E(\hat{\mathbf{g}}_i\hat{\mathbf{g}}_i')$) $\E(\hat{\mathbf{g}}_i\hat{\mathbf{g}}_i')$
  is nonsingular, $\hat{\mathbf{g}}_i \equiv \hat{\boldsymbol{\eta}}_i \otimes \mathbf{x}_i$.

  \textbf{Hint:} For the orthogonality conditions, note that \eqref{eq:5.1.8b} can be written as
  $\E(\boldsymbol{\eta}_i \otimes \mathbf{x}_i) = \mathbf{0}$ and that
  $\hat{\boldsymbol{\eta}}_i \otimes \mathbf{x}_i = (\mathbf{C}' \otimes \mathbf{I}_K)(\boldsymbol{\eta}_i \otimes \mathbf{x}_i)$.
  The identification
  condition is equivalent to \eqref{eq:5.1.15} (see the answer sheet for proof). To show
  that $\E(\hat{\boldsymbol{\eta}}_i\hat{\boldsymbol{\eta}}_i' \mid \mathbf{x}_i)$ does not depend on
  $\mathbf{x}_i$, use the fact that $\hat{\boldsymbol{\eta}}_i = \mathbf{C}'\boldsymbol{\eta}_i = \mathbf{C}'\boldsymbol{\varepsilon}_i$
  since $\boldsymbol{\varepsilon}_i = \mathbf{1}_M \cdot \alpha_i + \boldsymbol{\eta}_i$. Given conditional homoskedasticity, we have
  \[
    \E(\hat{\mathbf{g}}_i\hat{\mathbf{g}}_i')
    = \E(\hat{\boldsymbol{\eta}}_i\hat{\boldsymbol{\eta}}_i') \otimes \E(\mathbf{x}_i\mathbf{x}_i').
  \]
  So the last condition (nonsingularity of $\E(\hat{\mathbf{g}}_i\hat{\mathbf{g}}_i')$) is simply that
  $\E(\mathbf{x}_i\mathbf{x}_i')$ be nonsingular.

\item Show that the fixed-effects estimator \eqref{eq:5.2.4} can be written in the GMM
  format (7) with
  \begin{equation}
    \mathbf{S}_{\mathbf{xz}} = \frac{1}{n}\sum_{i=1}^{n}\hat{\mathbf{F}}_i \otimes \mathbf{x}_i,
    \quad
    \mathbf{s}_{\mathbf{xy}} = \frac{1}{n}\sum_{i=1}^{n}\hat{\mathbf{y}}_i \otimes \mathbf{x}_i,
    \quad
    \hat{\mathbf{W}} = (\mathbf{C}'\mathbf{C})^{-1} \otimes
    \biggl(\frac{1}{n}\sum_{i=1}^{n}\mathbf{x}_i\mathbf{x}_i'\biggr)^{\!-1}.
    \tag{10}
  \end{equation}

  \textbf{Hint:} Use the following result from matrix algebra: If $\mathbf{C}$ satisfies condition
  (8), then
  \begin{equation}
    \mathbf{C}(\mathbf{C}'\mathbf{C})^{-1}\mathbf{C}' = \mathbf{Q} \tag{11}
  \end{equation}
  (the annihilator associated with $\mathbf{1}_M$). The elements of $\mathbf{f}_{im}$ are a subset of
  the elements of $\mathbf{x}_i$. Observe the ``disappearance of $\mathbf{x}$.'' So, for example,
  \[
    \frac{1}{n}\sum_{i=1}^{n}\mathbf{f}_{im}\mathbf{x}_i'
    \biggl(\frac{1}{n}\sum_{i=1}^{n}\mathbf{x}_i\mathbf{x}_i'\biggr)^{\!-1}
    \frac{1}{n}\sum_{i=1}^{n}\mathbf{x}_i\mathbf{f}_{ih}'
    = \frac{1}{n}\sum_{i=1}^{n}\mathbf{f}_{im}\mathbf{f}_{ih}'.
  \]

\item Let $\hat{\boldsymbol{\Psi}}$ ($(M-1) \times (M-1)$) be a consistent estimate of
  $\boldsymbol{\Psi} \equiv \E(\hat{\boldsymbol{\eta}}_i\hat{\boldsymbol{\eta}}_i')$. Show
  that the efficient GMM estimator is given by
  \begin{equation}
    \hat{\boldsymbol{\beta}} = \biggl(\frac{1}{n}\sum_{i=1}^{n}\hat{\mathbf{F}}_i'
    \hat{\boldsymbol{\Psi}}^{-1}\hat{\mathbf{F}}_i\biggr)^{\!-1}
    \frac{1}{n}\sum_{i=1}^{n}\hat{\mathbf{F}}_i'\hat{\boldsymbol{\Psi}}^{-1}\hat{\mathbf{y}}_i. \tag{12}
  \end{equation}
  Show that
  \begin{equation}
    \avar(\hat{\boldsymbol{\beta}})
    = \bigl[\E(\hat{\mathbf{F}}_i'\boldsymbol{\Psi}^{-1}\hat{\mathbf{F}}_i)\bigr]^{-1}. \tag{13}
  \end{equation}
  \textbf{Hint:} Apply Proposition~4.7 to the present system of $M - 1$ equations.

\item Show that the following is consistent for $\boldsymbol{\Psi}$:
  \begin{equation}
    \underset{((M-1) \times (M-1))}{\hat{\boldsymbol{\Psi}}}
    = \frac{1}{n}\sum_{i=1}^{n}
    (\hat{\mathbf{y}}_i - \hat{\mathbf{F}}_i\hat{\boldsymbol{\delta}}_{\text{FE}})
    (\hat{\mathbf{y}}_i - \hat{\mathbf{F}}_i\hat{\boldsymbol{\delta}}_{\text{FE}})'. \tag{14}
  \end{equation}
  \textbf{Hint:} Verify that all the assumptions of Proposition~4.1 are satisfied.

\item (Sargan's statistic)\quad Derive Sargan's statistic associated with the efficient
  GMM estimator. \textbf{Hint:} In Proposition~4.7(c), set
  \begin{equation}
    \hat{\mathbf{S}} = \hat{\boldsymbol{\Psi}} \otimes
    \biggl(\frac{1}{n}\sum \mathbf{x}_i\mathbf{x}_i'\biggr), \quad
    \mathbf{g}_n(\hat{\boldsymbol{\beta}})
    = \mathbf{s}_{\mathbf{xy}} - \mathbf{S}_{\mathbf{xz}}\hat{\boldsymbol{\beta}}.
  \end{equation}

\item Suppose that $\boldsymbol{\eta}_i$ is spherical in that
  \begin{equation}
    \E(\boldsymbol{\eta}_i\boldsymbol{\eta}_i') = \sigma_\eta^2\,\mathbf{I}_M. \tag{15}
  \end{equation}
  Verify that, if $\hat{\sigma}_\eta^2$ is a consistent estimator of $\sigma_\eta^2$, then
  \begin{equation}
    \hat{\boldsymbol{\Psi}} = \hat{\sigma}_\eta^2\,\mathbf{C}'\mathbf{C} \tag{16}
  \end{equation}
  is consistent for $\boldsymbol{\Psi}$. Verify that the efficient GMM estimator (12) with this
  choice of $\hat{\boldsymbol{\Psi}}$ is numerically equal to the fixed-effects estimator.

\item (Invariance to the choice of $\mathbf{C}$)\quad Verify: if an $M \times (M-1)$ matrix $\mathbf{C}$ satisfies
  (8), so does $\mathbf{B} = \mathbf{C}\mathbf{A}$ where $\mathbf{A}$ is an $(M-1) \times (M-1)$ nonsingular matrix.
  Replace $\mathbf{C}$ everywhere in this section by $\mathbf{B}$ and verify that the choice of $\mathbf{C}$
  does not change any of the results of this question.
\end{enumerate}

\exercise[3]
Let $SSR$ be the sum of squared residuals defined in \eqref{eq:5.2.15}. For the error-components
model with the spherical assumption that
$\E(\boldsymbol{\eta}_i\boldsymbol{\eta}_i') = \sigma_\eta^2\,\mathbf{I}_M$, prove
\[
  \plim_{n \to \infty} \frac{SSR}{n} = (M - 1) \cdot \sigma_\eta^2.
\]

\textbf{Hint:}
\begin{align*}
  SSR &= \sum_{i=1}^{n}[\mathbf{Q}(\mathbf{y}_i - \mathbf{F}_i\hat{\boldsymbol{\beta}}_{\text{FE}})]'
  [\mathbf{Q}(\mathbf{y}_i - \mathbf{F}_i\hat{\boldsymbol{\beta}}_{\text{FE}})] \\
  &= \sum_{i=1}^{n}(\mathbf{y}_i - \mathbf{F}_i\hat{\boldsymbol{\beta}}_{\text{FE}})'\mathbf{Q}
  (\mathbf{y}_i - \mathbf{F}_i\hat{\boldsymbol{\beta}}_{\text{FE}}) \\
  &\quad (\text{since } \mathbf{Q} \text{ is symmetric and idempotent}) \\
  &= \sum_{i=1}^{n}(\mathbf{y}_i - \mathbf{F}_i\hat{\boldsymbol{\beta}}_{\text{FE}})'
  \mathbf{C}(\mathbf{C}'\mathbf{C})^{-1}\mathbf{C}'
  (\mathbf{y}_i - \mathbf{F}_i\hat{\boldsymbol{\beta}}_{\text{FE}})
  \quad (\text{by (11)}) \\
  &= \sum_{i=1}^{n}\hat{\mathbf{v}}_i'(\mathbf{C}'\mathbf{C})^{-1}\hat{\mathbf{v}}_i
  \quad (\hat{\mathbf{v}}_i \equiv \mathbf{C}'\mathbf{y}_i
  - \mathbf{C}'\mathbf{F}_i\hat{\boldsymbol{\beta}}_{\text{FE}}) \\
  &= \sum_{i=1}^{n}\operatorname{trace}[\hat{\mathbf{v}}_i'(\mathbf{C}'\mathbf{C})^{-1}\hat{\mathbf{v}}_i] \\
  &= \sum_{i=1}^{n}\operatorname{trace}[(\mathbf{C}'\mathbf{C})^{-1}\hat{\mathbf{v}}_i\hat{\mathbf{v}}_i'] \\
  &= n \cdot \operatorname{trace}\biggl[(\mathbf{C}'\mathbf{C})^{-1}
  \frac{1}{n}\sum_{i=1}^{n}\hat{\mathbf{v}}_i\hat{\mathbf{v}}_i'\biggr].
\end{align*}

With the usual technical condition (that $\E(\mathbf{f}_{im}\mathbf{f}_{ih}')$ exists and is finite for all $m$, $h$,
which obtains if $\E(\mathbf{x}_i\mathbf{x}_i')$ exists and is finite), it should now be routine to show
that
\[
  \frac{1}{n}\sum_{i=1}^{n}\hat{\mathbf{v}}_i\hat{\mathbf{v}}_i'
  \pto \E(\mathbf{v}_i\mathbf{v}_i') \quad \text{as } n \to \infty,
\]
where $\mathbf{v}_i \equiv \mathbf{C}'\mathbf{y}_i - \mathbf{C}'\mathbf{F}_i\boldsymbol{\beta}$.
Show that $\E(\mathbf{v}_i\mathbf{v}_i') = \sigma_\eta^2 \cdot \mathbf{C}'\mathbf{C}$.

\exercise[4]
(Inconsistency of the fixed-effects estimator for dynamic models)\quad
Consider a
``dynamic'' model
\[
  y_{im} = \alpha_i + \rho \cdot y_{i,m-1} + \eta_{im}
  \quad (m = 1, 2, \ldots, M;\; i = 1, 2, \ldots, n).
\]

Suppose we have a random sample on $(y_{i0}, y_{i1}, \ldots, y_{iM})$. We assume that
$\E(\alpha_i\,\eta_{im}) = 0$, and $\E(y_{i0}\,\eta_{im}) = 0$ for all $m$. Also assume there is no cross-equation
correlation: $\E(\eta_{im}\,\eta_{ih}) = 0$ for $m \neq h$, and that the second moment
across equations is the same: $\E(\eta_{im}^2) = \sigma_\eta^2$.

\begin{enumerate}[label=(\alph*)]
\item Write this system of $M$ equations as \eqref{eq:5.1.1pp} on page~329 by properly defining
  $\mathbf{y}_i$, $\mathbf{F}_i$, and $\mathbf{b}_i$.

\item Show: $\E(y_{im} \cdot \eta_{ih}) = 0$ if $m < h$. \textbf{Hint:}
  \[
    y_{im} = \eta_{im} + \rho\eta_{i,m-1} + \cdots + \rho^{m-1}\eta_{i1}
    + \frac{1 - \rho^m}{1 - \rho}\,\alpha_i + \rho^m y_{i0}.
  \]

\item Verify that $\E(y_{im} \cdot \eta_{ih}) = \rho^{m-h}\sigma_\eta^2$ for $m \geq h$.
  (optional)\quad Show that
  \[
    \E(\mathbf{F}_i'\mathbf{Q}\boldsymbol{\eta}_i)
    = -\frac{\sigma_\eta^2}{M} \cdot \frac{(M-1) - M\rho + \rho^M}{(1-\rho)^2}.
  \]

\item Therefore, if $\E(\mathbf{F}_i'\mathbf{Q}\mathbf{F}_i)$ is nonsingular and if
  $(M-1) - M\rho + \rho^M \neq 0$, then
  the fixed-effects estimator for this dynamic model is \emph{inconsistent}. Which
  assumption of Proposition~5.1 guaranteeing the consistency of the fixed-effects
  estimator is violated?
\end{enumerate}

\exercise[5]
(Optional, identification ensures nonsingularity)

\begin{enumerate}[label=(\alph*)]
\item Show that $\E(\widetilde{\mathbf{F}}_i'\widetilde{\mathbf{F}}_i)$ in \eqref{eq:5.2.7} is invertible under assumption \eqref{eq:5.1.15}, that
  $\E(\mathbf{Q}\mathbf{F}_i \otimes \mathbf{x}_i)$ is of full column rank. \textbf{Hint:}
  \[
    \E(\widetilde{\mathbf{F}}_i'\widetilde{\mathbf{F}}_i)
    = \E(\mathbf{F}_i'\mathbf{Q}\mathbf{F}_i)
    = \sum_{m=1}^{M}\sum_{h=1}^{M}q_{mh}\,\E(\mathbf{f}_{im}\mathbf{f}_{ih}')
  \]
  where $q_{mh} = (m,h)$ element of $\mathbf{Q}$. Show that
  \[
    \E(\mathbf{f}_{im}\mathbf{f}_{ih}')
    = \E(\mathbf{f}_{im}\mathbf{x}_i')\,[\E(\mathbf{x}_i\mathbf{x}_i')]^{-1}\,
      \E(\mathbf{x}_i\mathbf{f}_{ih}').
  \]
  (Recall that the elements of $\mathbf{f}_{im}$ are included in $\mathbf{x}_i$.) Finally, show that
  \begin{align*}
    \mathbf{x}_i\mathbf{f}_{im}' &= \mathbf{f}_{im}' \otimes \mathbf{x}_i,\;
    (\mathbf{F}_i \otimes \mathbf{x}_i)' = (\mathbf{f}_{i1} \otimes \mathbf{x}_i', \ldots,
    \mathbf{f}_{iM} \otimes \mathbf{x}_i'), \text{ and} \\
    \E(\widetilde{\mathbf{F}}_i'\widetilde{\mathbf{F}}_i)
    &= \E(\mathbf{F}_i \otimes \mathbf{x}_i)'\bigl[\mathbf{Q} \otimes [\E(\mathbf{x}_i\mathbf{x}_i')]^{-1}\bigr]
    \,\E(\mathbf{F}_i \otimes \mathbf{x}_i).
  \end{align*}

\item Use a similar argument to show that
  $\E\!\bigl[\widetilde{\mathbf{F}}_i'\,\E(\tilde{\boldsymbol{\eta}}_i\tilde{\boldsymbol{\eta}}_i')\widetilde{\mathbf{F}}_i\bigr]$
  in \eqref{eq:5.2.7} is nonsingular.
  \textbf{Hint:} Since
  \[
    \tilde{\boldsymbol{\eta}}_i = \mathbf{Q}\boldsymbol{\eta}_i = \mathbf{Q}\boldsymbol{\varepsilon}_i
  \]
  we have
  \[
    \widetilde{\mathbf{F}}_i'\,\E(\tilde{\boldsymbol{\eta}}_i\tilde{\boldsymbol{\eta}}_i')\widetilde{\mathbf{F}}_i
    = \widetilde{\mathbf{F}}_i'\mathbf{Q}\,\E(\boldsymbol{\varepsilon}_i\boldsymbol{\varepsilon}_i')\,
      \mathbf{Q}\widetilde{\mathbf{F}}_i
    = \widetilde{\mathbf{F}}_i'\,\E(\boldsymbol{\varepsilon}_i\boldsymbol{\varepsilon}_i')\,\widetilde{\mathbf{F}}_i.
  \]
  Also,
  \[
    \E\!\bigl[\widetilde{\mathbf{F}}_i'\,\E(\boldsymbol{\varepsilon}_i\boldsymbol{\varepsilon}_i')\widetilde{\mathbf{F}}_i\bigr]
    = \E(\mathbf{Q}\mathbf{F}_i \otimes \mathbf{x}_i)'\bigl(\E(\boldsymbol{\varepsilon}_i\boldsymbol{\varepsilon}_i')
    \otimes [\E(\mathbf{x}_i\mathbf{x}_i')]^{-1}\bigr)\,\E(\mathbf{Q}\mathbf{F}_i \otimes \mathbf{x}_i).
  \]

\item Show that $\E(\hat{\mathbf{F}}_i'\boldsymbol{\Psi}^{-1}\hat{\mathbf{F}}_i)$ in (13) of Analytical Exercise~2 is nonsingular.
  \textbf{Hint:} As shown in Analytical Exercise~2(c),
  $\E(\mathbf{Q}\mathbf{F}_i \otimes \mathbf{x}_i)$ is of full column rank
  if and only if $\E(\mathbf{C}'\mathbf{F}_i \otimes \mathbf{x}_i)$ is of full column rank.
\end{enumerate}

\exercise[6]
(Optional, alternative identification condition)\quad
Consider the error-component
model \eqref{eq:5.1.1pp} on page~329. Let $p \equiv \#\boldsymbol{\beta}$. If $\mathbf{A}_m$ is the matrix that picks up
$\mathbf{f}_{im}$ from $\mathbf{x}_i$, it can be written as
\[
  \underset{(K \times p)}{\mathbf{A}_m} = \begin{bmatrix} \mathbf{e}_m \otimes \mathbf{I}_p \\ \mathbf{0} \end{bmatrix},
\]
where $\mathbf{e}_m$ is the $m$-th row of $\mathbf{I}_M$.

\begin{enumerate}[label=(\alph*)]
\item Let $\mathbf{A}$ ($KM \times p$) be the matrix obtained from stacking
  $(\mathbf{A}_1, \ldots, \mathbf{A}_M)$. Show
  that $\mathbf{F}_i = (\mathbf{I}_M \otimes \mathbf{x}_i')\mathbf{A}$.

\item Show that the rank condition \eqref{eq:5.1.15}, combined with the condition that
  $\E(\mathbf{x}_i\mathbf{x}_i')$ be nonsingular (which is implied by \eqref{eq:5.1.5} and
  \eqref{eq:5.1.6}), can be
  written as
  \[
    \operatorname{rank}[(\mathbf{Q} \otimes \mathbf{I}_K)\mathbf{A}] = p
    \quad (\equiv \#\boldsymbol{\beta}).
  \]

  \textbf{Hint:}
  \begin{align*}
    \mathbf{F}_i \otimes \mathbf{x}_i
    &= ([\mathbf{I}_M \otimes \mathbf{x}_i')\mathbf{A}] \otimes \mathbf{x}_i \\
    &= [(\mathbf{I}_M \otimes \mathbf{x}_i') \otimes \mathbf{x}_i](\mathbf{A} \otimes 1) \\
    &= (\mathbf{I}_M \otimes \mathbf{x}_i' \otimes \mathbf{x}_i)\mathbf{A} \\
    &= [\mathbf{I}_M \otimes (\mathbf{x}_i\mathbf{x}_i')]\mathbf{A}.
  \end{align*}
  So
  \[
    \widetilde{\mathbf{F}}_i \otimes \mathbf{x}_i
    = (\mathbf{Q} \otimes \mathbf{I}_K)(\mathbf{F}_i \otimes \mathbf{x}_i)
    = [\mathbf{Q} \otimes (\mathbf{x}_i\mathbf{x}_i')]\mathbf{A}.
  \]
  So $\E(\widetilde{\mathbf{F}}_i \otimes \mathbf{x}_i)
  = [\mathbf{I}_M \otimes \E(\mathbf{x}_i\mathbf{x}_i')](\mathbf{Q} \otimes \mathbf{I}_K)\mathbf{A}$.
  (This result will be used in the next exercise.)
\end{enumerate}

\exercise[7]
(Optional, checking nonsingularity)\quad Prove that
$\E(\widetilde{\mathbf{F}}_i'\tilde{\boldsymbol{\eta}}_i\tilde{\boldsymbol{\eta}}_i'\widetilde{\mathbf{F}}_i)$
in \eqref{eq:5.2.23} is nonsingular, by taking the following steps.

\begin{enumerate}[label=(\alph*)]
\item Write $\widetilde{\mathbf{F}}_i'\tilde{\boldsymbol{\eta}}_i$ as a linear function of
  $\mathbf{g}_i \equiv \boldsymbol{\varepsilon}_i \otimes \mathbf{x}_i$.
  \textbf{Hint:} Let $\mathbf{A}$ be as in the
  previous exercise.
  \begin{align*}
    \widetilde{\mathbf{F}}_i'\tilde{\boldsymbol{\eta}}_i
    &= \mathbf{A}'(\mathbf{I}_M \otimes \mathbf{x}_i)\mathbf{Q}\boldsymbol{\eta}_i \\
    &= \mathbf{A}'(\mathbf{I}_M \otimes \mathbf{x}_i)\mathbf{Q}\boldsymbol{\varepsilon}_i
    \quad (\text{since } \boldsymbol{\varepsilon}_i = \mathbf{1}_M \cdot \alpha_i + \boldsymbol{\eta}_i) \\
    &= \mathbf{A}'(\mathbf{I}_M \otimes \mathbf{x}_i)(\mathbf{Q}\boldsymbol{\varepsilon}_i \otimes 1) \\
    &= \mathbf{A}'(\mathbf{Q}\boldsymbol{\varepsilon}_i \otimes \mathbf{x}_i) \\
    &= \mathbf{A}'(\mathbf{Q} \otimes \mathbf{I}_K)\mathbf{g}_i.
  \end{align*}

\item Prove the desired result. \textbf{Hint:} $\E(\mathbf{g}_i\mathbf{g}_i')$ is positive definite by
  \eqref{eq:5.1.6}.
\end{enumerate}

\subsection*{Empirical Exercises}

In this exercise, you will estimate the speed of convergence using several different
estimation techniques. We use the international panel constructed by Summers and
Heston (1991), which has become the standard data set for studying the growth of
nations. The Summers-Heston panel (also called the Penn World Table) includes
major macro variables measured on a consistent basis for virtually all the countries
of the world.

The file \texttt{SUM\_HES.ASC} is an extract for 1960--1985 from version~5.6 of the
Penn World Table downloaded from the NBER's home page (\texttt{www.nber.org}). The
file is organized as follows. For each country, information on the country is contained
in multiple records (rows). The first record has the country's identification
code and two dummies, one for a communist regime (referred to as \emph{COM}
below) and the other for the Organization of Petroleum Exporting Countries (referred to as
\emph{OPEC} below). The second record has the year, the population (in thousands), real
GDP per capita (in 1985 U.S.\ dollars), and the saving rate (in percent) for 1960.
The third record has the same for 1961, and so forth. The twenty-seventh record of
the country is for 1985.

The file contains all 125 countries for which the data are available. (The mapping
between the country ID and the country name is in \texttt{country.asc}. Country~1
is Algeria, 2 is Angola, etc. Note that the mapping in Penn World Table version
5.6 is different from that in Summers and Heston (1991).) GDP per capita is the
country's real GDP converted into U.S.\ dollars using the purchasing power parity
of the country's currency against the dollar for 1985 (the variable ``RGDPCH'' in
the Penn World Table). The saving rate is measured as the ratio of real investment
to real GDP (``I'' in the Penn World Table). See Sections~II and III (particularly
III.D) of Summers and Heston (1991) for how the variables are constructed.

The first issue we examine empirically is conditional convergence. (At this
juncture it would be useful to skim Mankiw, Romer, and Weil (1992).) Let $s_i$ be
the saving rate of country $i$ and $n_i$ be the country's population growth rate (which
we take to be the growth rate of labor). In the Solow-Swan growth model with
a Cobb-Douglas production function, the steady-state level of output per effective
labor, denoted $q^*$ in the text, is a linear function of
$\log(s_i) - \log(n_i + g + \delta)$, where
$g$ is the rate of labor-augmenting technical progress and $\delta$ is the depreciation rate
of the capital stock. Then, assuming $A(0)$, the initial level of technology, to be the
same across countries, equation \eqref{eq:5.4.10} can be written as
\begin{equation}
  y_{im} = \phi_m + \rho\, y_{i,m-1} + \gamma_1 \log(s_i)
  + \gamma_2 \log(n_i + g + \delta) + \eta_{im}. \tag{1}
\end{equation}

The growth model implies that $\gamma_1 = -\gamma_2$, but we will not impose this restriction.
As in Mankiw et al., assume $g + \delta$ to be 5 percent for all countries, and take $s_i$
to be the saving rate averaged over the 1960--1985 period. We measure $n_i$ as the
average annual population growth over the same period (Mankiw et al.\ uses the
average growth rate of the working-age population). Because our sample of 125
countries includes (former) communist countries and OPEC, we add
$COM$ ($=1$ if communist, 0 otherwise) and $OPEC$ ($=1$ if OPEC, 0 otherwise) to the equation:
\begin{equation}
  y_{im} = \phi_m + \rho\, y_{i,m-1} + \gamma_1 \log(s_i)
  + \gamma_2 \log(n_i + g + \delta)
  + \gamma_3 COM_i + \gamma_4 OPEC_i + \eta_{im}. \tag{2}
\end{equation}

\begin{enumerate}[label=(\alph*)]
\item By subtracting $y_{i,m-1}$ from both sides, we can rewrite (2) as
  \begin{align}
    y_{im} - y_{i,m-1} &= \phi_m + (\rho - 1)y_{i,m-1}
    + \gamma_1 \log(s_i) + \gamma_2 \log(n_i + g + \delta)
    \notag \\
    &\quad + \gamma_3 COM_i + \gamma_4 OPEC_i + \eta_{im}. \tag{2$'$}
  \end{align}

  For this specification, set $M = 1$ (just one equation to estimate) and take $t_0 = 1960$
  and $t_1 = 1985$, so the dependent variable, $y_{i1} - y_{i0}$, is the cumulative
  growth between 1960 and 1985.
  Plot $y_{i1} - y_{i0}$ against $y_{i0}$ for the 125 countries
  included in the sample. Is there any relation between the initial GDP ($y_{i0}$)
  and the subsequent growth ($y_{i1} - y_{i0}$)? Assuming outright that the error term is
  orthogonal to the regressors, estimate equation (2$'$) by OLS. Calculate the value
  of $\lambda$ (speed of convergence) implied by the OLS estimate of $\rho$. (It should be
  less than 1 percent per year.) Calculate the asymptotic standard error (i.e., $1/n$
  times the square root of the asymptotic variance) of your estimate of $\lambda$.
  \textbf{Hint:}
  By \eqref{eq:5.4.3}, $\hat{\lambda} = -\log(\hat{\rho})/25$. By the delta method (Lemma~2.5 of Section~2.1),
  \[
    \avar(\hat{\lambda}) = \frac{\avar(\hat{\rho})}{(25\hat{\rho})^2}.
  \]
  So the standard error of $\hat{\lambda} = (\text{standard error of } \hat{\rho})/(25\hat{\rho})$.
  Can you confirm the
  finding that ``if countries did not vary in their investment and population growth
  rates, there would be a strong tendency for poor countries to grow faster than
  rich ones'' (Mankiw et al., 1992, p.~428).

\item (Fixed-effects estimation)\quad
  By setting $M = 25$, $t_0 = 1960$,
  $t_1 = 1961, \ldots, t_{25} = 1985$, we can form a system of 25 equations, with
  \eqref{eq:5.4.10} as the
  $m$-th equation. Estimate the system by the fixed-effect technique, assuming
  that the error is spherical in the sense of \eqref{eq:5.2.12}. What is the implied value
  of $\lambda$? (It should be about 6.4 percent per year.) (Optional: Calculate the standard
  errors without the spherical error assumption. As emphasized in the text,
  the fixed-effects estimator is \emph{not} consistent. We nevertheless apply blindly the
  fixed-effect technique just to gain programming experience.)

  \textbf{TSP Tip:} TSP's command for the fixed-effects estimator (with the spherical
  error assumption) is \texttt{PANEL}. It, however, does not accept the data organization
  of \texttt{SUM\_HES.ASC}. It requires that the group (country) ID and the year
  be included in the rows containing information on group-years. The data file
  that \texttt{PANEL} can accept is \texttt{FIXED.ASC}. Its first few records are:

  \begin{verbatim}
1 0 0  1961  1599  1723
1 0 0  1962  1275  1599
     :
1 0 0  1984  2962  2903
1 0 0  1985  2988  2962
2 0 0  1961   976   931
2 0 0  1962   984   976
     :
  \end{verbatim}

  Here, the first variable is the country ID, followed by \emph{COM}, \emph{OPEC}, year, real
  per capita GDP for the year, and real per capita GDP for the previous year.

\item (optional)\quad As in the previous question, take $M = 25$ and $t_0 = 1960$,
  $t_1 = 1961, \ldots, t_{25} = 1985$. Then \eqref{eq:5.4.11} is a system of 24 equations if
  written as the model of Section~4.6:
  \begin{equation}
    y_{im} = \mathbf{z}_{im}'\boldsymbol{\delta} + \boldsymbol{\varepsilon}_{im}
    \quad (i = 1, 2, \ldots, 125;\; m = 1, 2, \ldots, 24), \tag{3}
  \end{equation}
  where the $y_{im}$ in (3) equals $y_{i,m+1} - y_{im}$ in \eqref{eq:5.4.11},
  \[
    \mathbf{z}_{im} = (0, \ldots, 0, 1, 0, \ldots, 0, y_{im} - y_{i,m-1})',
  \]
  a 25-dimensional vector whose $m$-th element is 1,
  \[
    \boldsymbol{\delta} = (\mu_1, \ldots, \mu_{24}, \rho)'.
  \]

  If $s_{im}$ is country $i$'s saving rate (investment/GDP ratio) in year $m$, use a constant
  and $s_{i,m-1}$ as the instrument in the $m$-th equation (so for example, in the
  first equation where the nonconstant regressor is $y_{i,1961} - y_{i,1960}$, the vector of
  instruments is $(1, s_{i,1960})')$.
  Apply the multiple-equation GMM technique to (3)
  assuming conditional homoskedasticity. (It will be very clumsy to use TSP
  or RATS; use GAUSS.) The estimator is (4.6.6) with $\hat{\mathbf{W}}$ set to the inverse of
  (4.5.3).) The implied value of $\lambda$ should be about 36 percent per year!
\end{enumerate}

\end{problemset}

\begin{answers}

\subsection*{Analytical Exercises}

\answer{1c}
$\E(\boldsymbol{\eta}_i \mid \mathbf{W})$

$= \E(\boldsymbol{\eta}_i \mid \mathbf{F})$ \quad (since $\mathbf{D}$ is a matrix of constants)

$= \E(\boldsymbol{\eta}_i \mid \mathbf{F}_1, \ldots, \mathbf{F}_n)$

$= \E(\boldsymbol{\eta}_i \mid \mathbf{F}_i)$

\quad (since $\boldsymbol{\eta}_i$ is indep.\ of
$(\mathbf{F}_1, \ldots, \mathbf{F}_{i-1}, \mathbf{F}_{i+1}, \ldots, \mathbf{F}_n)$
by i.i.d.\ assumption)

$= \mathbf{0}$.

Therefore, the regressors are exogenous. Also,
\begin{align*}
  \E(\boldsymbol{\eta}_i\boldsymbol{\eta}_i' \mid \mathbf{W})
  &= \E(\boldsymbol{\eta}_i\boldsymbol{\eta}_i' \mid \mathbf{F}) \\
  &= \E(\boldsymbol{\eta}_i\boldsymbol{\eta}_i' \mid \mathbf{F}_i) \\
  &= \sigma_\eta^2\,\mathbf{I}_M.
\end{align*}

By the i.i.d.\ assumption,
$\E(\boldsymbol{\eta}_i\boldsymbol{\eta}_j' \mid \mathbf{W}) = \mathbf{0}$ for $i \neq j$.
So $\E(\boldsymbol{\eta}\boldsymbol{\eta}' \mid \mathbf{W}) = \sigma_\eta^2\,\mathbf{I}_{Mn}$.

\medskip
\answer{2c}
Proof that ``$\E(\mathbf{Q}\mathbf{F}_i \otimes \mathbf{x}_i)$ is of full column rank''
$\Leftrightarrow$
``$\E(\mathbf{C}'\mathbf{F}_i \otimes \mathbf{x}_i)$ is of full column rank'':

Let $\mathbf{A} \equiv \E(\mathbf{Q}\mathbf{F}_i \otimes \mathbf{x}_i)$ and
$\mathbf{B} \equiv \E(\mathbf{C}'\mathbf{F}_i \otimes \mathbf{x}_i)$.
There exists an $M \times (M-1)$
matrix $\mathbf{L}$ of full column rank such that $\mathbf{Q} = \mathbf{L}\mathbf{C}'$
(it is actually $\mathbf{C}(\mathbf{C}'\mathbf{C})^{-1}$ by
(11)). Let $\mathbf{D} \equiv \mathbf{L} \otimes \mathbf{I}_K$. Then $\mathbf{A} = \mathbf{D}\mathbf{B}$.
Since the rank of a product of two
matrices is less than or equal to the rank of either of the two matrices,
\begin{equation}
  \operatorname{rank}(\mathbf{A}) \leq \operatorname{rank}(\mathbf{B}). \tag{1}
\end{equation}

Next consider a product $\mathbf{D}'\mathbf{A} = \mathbf{D}'\mathbf{D}\mathbf{B}$.
\begin{equation}
  \operatorname{rank}(\mathbf{D}'\mathbf{D}\mathbf{B})
  \leq \operatorname{rank}(\mathbf{D}'\mathbf{A})
  \leq \operatorname{rank}(\mathbf{A}). \tag{2}
\end{equation}

But since $\mathbf{L}$ is of full column rank, $\mathbf{D}'\mathbf{D}$ is nonsingular. Since multiplication by
a nonsingular matrix does not alter rank,
\[
  \operatorname{rank}(\mathbf{D}'\mathbf{D}\mathbf{B}) = \operatorname{rank}(\mathbf{B}).
\]

From (1)--(3) it follows that $\operatorname{rank}(\mathbf{A}) = \operatorname{rank}(\mathbf{B})$.
Since $\mathbf{A}$ and $\mathbf{B}$ have the same
number of columns, the desired conclusion follows.

\medskip
\answer{4a}
\[
  \mathbf{y}_i = \begin{bmatrix} y_{i1} \\ \vdots \\ y_{iM} \end{bmatrix}, \quad
  \mathbf{F}_i = \begin{bmatrix} y_{i0} \\ \vdots \\ y_{i,M-1} \end{bmatrix}.
\]

There is no $\mathbf{b}_i$.

\medskip
\answer{4d}
Since $\E(\eta_{ih} \cdot y_{im}) \neq 0$ for $m \geq h$, some of the ``cross'' orthogonalities are not
satisfied.
\end{answers}
